\subsection{Пересечения исходных задач расширенной выборки}
Дополнительный аудит, описанный во время генерации Luna, охватывает все 1\,140 исходных задач и фиксированную выборку из 1\,000 задач. Он не читает сгенерированные программы или записи исходов. После нормализации пробелов все выбранные инструкции различаются. Точное сравнение синтаксических деревьев эталонных программ удаляет строки документации, сохраняя имена, литералы и структуру. Дополнительная проверка токенов использует явно описанную адаптацию метода поиска близких дубликатов~\cite{allamanis2019duplication}; сходство инструкций измеряется по последовательностям из пяти слов после удаления абзаца со стартовым кодом.

Объединение связей по сходству инструкций Жаккара не ниже 0,70, близости кода и точным совпадениям даёт 992 связные группы среди выбранных задач: семь групп содержат 15 задач, крупнейшая --- три. Пороги сходства только инструкций 0,50, 0,70 и 0,90 дают 996, 999 и 1\,000 групп. Это разбиения для проверки чувствительности, а не оценки числа независимых по смыслу задач. В частности, /130 и /131 имеют одинаковое синтаксическое дерево эталона. Между выборками также есть пересечения: инструкция и дерево эталона /1120 совпадают с /1121 из отдельного опыта на 80 задачах, а дерево эталона /826 --- с задачей разработки /814. Разные идентификаторы поэтому не гарантируют отсутствия пересечений исходного материала между исследованиями.

Все назначения сохраняются. Дополнительный анализ будет пересэмплировать эти группы, сохраняя среднее по задачам и три повтора внутри каждой задачи. Полученные условные интервалы будут отделены от зарегистрированных тестов; неодинаковые размеры групп и неизвестные зависимости остаются ограничениями~\cite{mackinnon2023clusters}. Полные отпечатки, связи и разбиения сохранены в \path{reproducibility/task_dependence/source-audit-v2}. Аудит исходных задач завершён; интервалы на основе исходов ожидают завершения серии Luna.
